\paragraph{The initial symplectic labels and flags.}
\label{par:count-two-labels}
We now change endpoints. At the rational target we again start in the marked
eigenbasis, so that later changes of lattice can be read on the same root spaces.

\begin{lemma}[The rational flag in the marked eigenbasis]
\label{lem:rational-rank-four-marking}
For the rational datum, there is an integral symplectic
eigenbasis diagonalizing the normalizing element \(\tau\).
In this basis the ordinary flag admits a lower unipotent
symplectic frame \(N\). Denote its four independent lower entries by
\(N_{10},N_{20},N_{30},N_{21}\). Let \(S_K\) consist of the degrees
of nonzero off-diagonal entries of the global matrices \(r(g)\),
\(g\in G\), over the fraction field \(K\). Write \(S^{(0)}\)
for the degrees of the global and independent frame entries
that are units on this initial lattice. Then
\[
 -w_\Psi\in S_K,\qquad
 S_K\cap\{w_\beta,-w_\alpha,-w_\vartheta\}\ne\varnothing,
 \qquad S^{(0)}\subseteq\{-w_\Psi\}.
\]
The initial residual global matrices are diagonal. Only a frame
entry can contribute to \(S^{(0)}\).
\end{lemma}

\begin{proof}
\proofstep{Place the ordinary flag in a chart with unit leading minors.}
For the rational point use the residual order
\((\bar\vartheta,\bar\Psi,\bar\alpha,1)\).
The designated Frobenius at \(\ell\) separates the
characters of the two spaces on the indices
\(\{0,2\}\) and \(\{1,3\}\). The residue of the
ordinary line has its character from the first
space and lies in that space. The two-dimensional
step of the Lagrangian flag splits across the two
spaces, by the Frobenius projectors, with a line in
each. The two spaces are paired by the alternating
form. If the first line has nonzero coordinate zero,
the Lagrangian plane projects isomorphically to the
first two coordinates. Indeed if its other line had
coordinate one zero it would be spanned by the
vector of index three, which pairs nontrivially
with the first line.

If instead coordinate zero vanishes, choose the signed symplectic permutation inducing the exchanges
\[
 (0\ 2)(1\ 3).
\]
This is allowed because the original first line
then has character both \(\bar\vartheta\) and
\(\bar\alpha\) on the marked decomposition group.
Their equality says \(\bar\Psi=1\) on the whole group,
and the other two characters coincide there as well.
The permutation acts on the global basis, with the
order of the given local inertia characters retained.
In either order the flag has unit first two leading
minors. The third is a unit by the relation between
the inverse matrix and its symplectic adjoint.

\proofstep{Diagonalize the marked global matrix without changing the residual flag.}
Let \(\lambda_i\) be the four Hensel lifts of the distinct
eigenvalues of the matrix \(R(\tau)\) in the complete
coefficient ring. The projectors
\[
 e_i=\prod_{j\ne i}\frac{R(\tau)-\lambda_j I}
                              {\lambda_i-\lambda_j}
\]
are integral because every denominator is a unit.
Applying \(e_i\) to the corresponding standard basis vector
gives an eigenvector reducing to that vector. For the
alternating form \(\langle\, , \,\rangle\) and the similitude
\(s(\tau)\), the eigenvector identity is
\[
  (\lambda_i\lambda_j-s(\tau))\langle v_i,v_j\rangle=0 .
\]
Only the paired indices \(0,3\) and \(1,2\) can have a nonzero
pairing, by their distinct residual eigenvalues. Each of
those pairings is a unit, congruent to the old pairing.
Rescale one vector of each pair to recover the symplectic
form with a change of basis congruent to the identity.

In this basis the unit leading minors give Gaussian
factorization of the flag frame as \(N U\), with \(N\) lower
unipotent and \(U\) upper triangular. Write \(J\) for the antidiagonal matrix of the form and \(\mu\) for the similitude of the frame \(A=NU\). The group automorphism \(\iota(X)=J^{-1}X^{-t}J\) preserves the lower unipotent and the upper triangular subgroups.

Since \(A^t J A=\mu J\), it gives a second factorization:
\[
 NU=A=\mu\iota(A)=\iota(N)\bigl(\mu\iota(U)\bigr).
\]
Consequently \(N^{-1}\iota(N)=U(\mu\iota(U))^{-1}\) is both lower unipotent and upper triangular. It is the identity, so \(\iota(N)=N\), equivalently \(N^tJN=J\). The upper factor preserves the ordinary flag
and its characters. Modulo the uniformizer the first two
columns follow their separate Frobenius spaces, and the
perpendicular conditions determine the two later steps.
Thus the only possible residual entries of \(N-I\) have
degree \(-w_\Psi\).

\proofstep{Locate the needed roots over the fraction field.}
If the rational representation is irreducible, the
algebra spanned by its matrices has all matrix-unit
positions in this eigenbasis. Let \(P_i\) be the diagonal projectors in this eigenbasis.
An arrow \(j\to i\) means that some action matrix \(r(g)\)
has \(r(g)_{ij}\ne0\). Then
\[
 P_i r(g)P_j=r(g)_{ij}E_{ij}.
\]
A proper terminal component of the directed graph would span
a proper stable coordinate subspace, so irreducibility gives
a path between each ordered pair of vertices. Repeated multiplication along a path removes each intermediate index:
\[
 E_{i_k i_{k-1}}\cdots E_{i_1 i_0}=E_{i_k i_0}.
\]
The projectors are polynomials in \(r(\tau)\). Hence every
product just obtained is in the \(K\)-linear span of the \(r(g)\):
products of group matrices are again group matrices. A nonzero
\((i,j)\)-coordinate in this span forces \(r(g)_{ij}\ne0\)
for some \(g\), with no assertion about its valuation.

In the block-extension case, the original matrices
preserve the space on indices \(1,3\), and the
diagonalizing eigenvectors in that block lie in the
same space. The change of basis is block triangular
and this block still spans the full matrix algebra.
On this block the matrix-unit calculation gives
\[
    (1,3),(3,1)\text{ in the matrix support over }K
       \quad\Longrightarrow\quad \{w_\Psi,-w_\Psi\}\subset S_K .
\]
 The nonsplit extension has a
nonzero cross-block entry. Otherwise the other
coordinate block gives a stable complement and
conjugating it back gives a stable complement
which is the graph of a map to the first block.
Such a graph is exactly the split extension
excluded in the rational datum. With rows \(1,3\) and columns \(0,2\), the four degrees are
\[
  \deg(1,2)=w_\beta,\quad \deg(3,0)=-w_\vartheta,\quad
  \deg(1,0)=\deg(3,2)=-w_\alpha .
\]

These entries give the two assertions about \(S_K\).
The initial global reduction is the given diagonal model and the
change by Hensel projectors is the identity on its reduction,
up to the preceding simultaneous permutation. The flag calculation
gives the asserted inclusion for \(S^{(0)}\).
\label{cor:cohomology-count-two-normalization-1}
\end{proof}

\paragraph{The torus normalization of the target.}
\label{par:count-two-normalization}
The rational eigenbasis fixes the labels but leaves a diagonal freedom. We use it to
reduce the target lattice to the shapes needed for interpolation.

\begin{proposition}[The two rational lattice shapes]
\label{prop:rational-rank-four-torus}
Assume the odd-root localization vanishing for the marked rational
datum and the basis of Lemma~\ref{lem:rational-rank-four-marking}.
There is a diagonal symplectic similitude, after adjoining
a square root of the uniformizer, giving an integral model
whose reduced global matrices \(r\) and reduced lower frame \(N\)
satisfy the following bounds. Here \(S\) counts the degrees of
nonzero reduced global entries off the diagonal and of the nonzero entries among the four
independent frame coordinates after the change of lattice:

\[
 S\subseteq\{w_\beta,-w_\Psi,-w_\alpha,-w_\vartheta\},
       \qquad r_{12}\not\equiv0,\qquad N_{21}=0.
\] Either degree \(-w_\Psi\) occurs in the residual
global matrices or it occurs only through \(N_{20}\ne0\),
in which case \(r_{12}\) cannot have the form
\(c(d_2-d_1)\) on the marked decomposition group.
\end{proposition}

\begin{proof}
\proofstep{Maximize the two torus coordinates subject to integrality.}
Use the normalized rational datum above in its small
coefficient model over a complete
discrete valuation ring with uniformizer \(\varpi\).
Use additive valuation normalized by
\(v(\varpi)=1\). Retain the four independent lower frame entries
\((1,0),(2,0),(3,0),(2,1)\); the symplectic equations
express the other lower entries in these. For each degree \(w=(w_1,w_2)\) occurring in a nonzero global or retained frame entry, define the minimum
\[
 k_w=\min\{v(a):a\ne0\text{ a global or retained frame entry of degree }w\}.
\] The given integral model makes
this a minimum of nonnegative integers. All global
off-diagonal entries, and the frame entries other than the
possible degree \(-w_\Psi\), have positive valuation.

Conjugation by the proposed diagonal with weights
\((a+b,b,a,0)\) changes a degree-\(w\) coefficient of
valuation \(k\) into one of valuation \(k+(w_1a+w_2b)/2\).
Thus integrality of every entry to be checked is equivalent
to requiring
\[
             w_1a+w_2b+2k_w\ge0
                    \qquad\text{for each degree }w\text{ with }k_w<\infty .
\]
The initial lattice makes \((0,0)\) feasible. Put \(k_w=+\infty\)
if there is no entry of degree \(w\). The two assertions about
\(S_K\) in Lemma~\ref{lem:rational-rank-four-marking} give
\[
 b\le 2k_{-w_\Psi},\qquad
 a\le e b+2k_{(-1,e)}\quad\text{for some }e\in\{-1,0,1\}.
\]
Each right side for its indicated degree is finite. Thus the feasible
integer pairs attain a maximum in \(b\), and on that slice in \(a\).

Take \(p^2=\varpi\) and use
the diagonal symplectic similitude with entries
\((p^{a+b},p^b,p^a,1)\). Subscripts of the
four ordered positions correspond to
\(((1,1),(0,1),(1,0),(0,0))\).

For the diagonal matrix \(D\) just defined, let \(w\) be the difference
of the two position degrees at an entry \(x_{ij}\). With \(v_\varpi(\varpi)=1\),
\[
 (DxD^{-1})_{ij}=p^{w_1a+w_2b}x_{ij},\qquad
 v_\varpi((DxD^{-1})_{ij})=v_\varpi(x_{ij})+(w_1a+w_2b)/2.
\]
The defining inequalities are exactly nonnegativity of these valuations
for the matrix and frame coordinates. For an original unit \(u\) and
a nonzero coordinate \(x=u\varpi^h\), this gives the integral element
\[
           p^{w_1a+w_2b}x=u p^{2h+w_1a+w_2b}
                     \in\mathcal O_{K(p)}.
\] Conjugation
preserves the diagonal, the ordinary flag equations
and the normalizing matrix at \(\tau\). This also
preserves the small-model condition and continuity.
Multiplication by finitely many fixed powers of \(p\)
only shifts the adic indices by fixed amounts.

\proofstep{Find the roots that prevent an increase in the maximizer.}
Reduction after conjugation gives the exact identity
\[
 S=\{w:k_w<\infty,\ w_1a+w_2b+2k_w=0\}.
\]
Indeed the minimum is attained in the discrete valuation ring,
and precisely entries of new valuation zero survive. Call these
degrees tight. In all subsequent root and flag equations, \(r\)
and \(N\) denote the reductions on this new lattice. For any tight root other
than \(-w_\Psi\), the value of \(w_1a+w_2b\)
is strictly negative. For \(-w_\Psi\) tight it gives
\(b\ge0\). The failed moves \((1,0)\) and \((0,1)\)
at the maximum require a tight root in each of
\[
 \{-w_\alpha,w_\beta,-w_\vartheta\},\qquad
 \{-w_\Psi,-w_\beta,-w_\vartheta\}.
\]
For either move the slack \(w_1a+w_2b+2k_w\) decreases
by at most one, so a failed move must meet a tight constraint.

At a tight root the integer \(w_1a+w_2b\) is even. The two
diagonal weights have linear forms \(a-b\) and \(a+b\)
up to sign. Tightness of either makes \(a+b\) even.
If neither blocking root is diagonal, the first gives
\(a\equiv0\pmod2\) and the second gives \(b\equiv0\pmod2\).
In either case the slack of a diagonal constraint is even. The failed moves \((-1,1)\) and \((1,1)\)
therefore require a tight root in each of
\[
 \{-w_\Psi,w_\alpha,-w_\beta\},\qquad
 \{-w_\Psi,-w_\alpha,-w_\vartheta\}.
\]
The only inequalities that drop by two in these moves
are diagonal roots, for which the preceding slack is
even. Hence slack one cannot block one of these moves.

For brevity put \(A=w_\alpha\), \(B=w_\beta\),
\(P=w_\Psi\), \(T=w_\vartheta\) in the following
table. The tight set is contained in one of the
listed sets and contains the two marked roots in its
row.
\[
\begin{array}{c|l}
\text{two tight roots}&\text{containing set}\\ \hline
B,-P & \{B,-P,-A,-T\}\\
-A,-P & \{-A,-P,-T\}\\
-T,-P & \{-T,-P\}\\
-A,-B & \{-A,-B,-P,-T\}\\
-T,-B & \{-T,-B,-P\}\\
-T,A & \{-T,A,-B,-P\}
\end{array}
\]
To check the table, suppose first that \(B\) is
tight, so \(b<a\) and \(-B\) is excluded.
The third failed move forces \(-P\) or \(A\).
If \(-P\) is absent, \(A\) gives \(b<a<0\), excluding
every root for the second move. Thus \(-P\) is
tight and \(a>b\ge0\), giving the first row.
If \(B\) is absent and \(-B\) is present, the first
move forces \(-A\) or \(-T\).

With \(-A\) tight
one has \(b>a>0\), giving row four. Otherwise the inequalities supplied by \(-T\) and \(-B\) give
\[
 0<a+b<2b\quad\Longrightarrow\quad b>0,
\]
where the strict upper bound uses \(a<b\).
According as \(A\) is absent or present one obtains
row five or six. If neither \(B\) nor \(-B\) is tight,
the second and third moves require either \(-P\) or
both \(-T,A\). These last two exclude \(-A,T,P\)
by their signs, giving row six. With \(-P\) tight,
\(b\ge0\) and the first move requires \(-A\) or \(-T\).
If \(-A\) is present, the signs exclude \(A,T,P\),
giving row two. In the case where \(-P\) is tight and \(-A\) is absent, \(-T\) must be present. The remaining test reads
\[
 A\text{ present}\ \Longrightarrow\ \text{row six},\qquad A\text{ absent}\ \Longrightarrow\ \text{row three}.
\]

\proofstep{Use localization to exclude unsupported root cocycles.}
Use the odd-root localization vanishing for this marked
rational datum from the rank-four cohomology calculation.
If \(r_{ij}\) has degree in \(\{\pm A,\pm B,\pm T\}\) and
only the diagonal terms in its multiplication law, then
\(c_{ij}=r_{ij}/d_j\) satisfies
\[
    c_{ij}(gh)=c_{ij}(g)+
                  (d_i/d_j)(g)c_{ij}(h).
\]
If its local restriction is a coboundary, the localization
statement makes it a global coboundary. In the marked
eigenbasis, \(c_{ij}(\tau)=0\), while \(d_i(\tau)/d_j(\tau)\ne1\).
This forces the global coboundary constant to be zero.
The same reasoning on the designated local Frobenius
forces a local constant to vanish whenever its root
character is nontrivial there.

If \(-A\) is absent, \(r_{10}=r_{32}=0\). Hence the two
intermediate terms of the \((3,0)\) entry vanish:
\[
 r_{31}(g)r_{10}(h)+r_{32}(g)r_{20}(h)=0.
\]
In the first-column flag equation the intermediate term involving \(N_{10}\) vanishes. The remaining identity is
\[
 r_{30}=(d_0-d_3)N_{30}\qquad(N_{10}=0).
\]
Thus the odd-root argument kills \(r_{30}\) and \(N_{30}\)
and excludes \(-T\) from the tight set.

If both \(B,P\)
are absent, \(r_{12}=r_{13}=0\); similarly the \((1,0)\)
entry has no intermediate terms, and its local equation
is \(r_{10}=(d_0-d_1)N_{10}\). Hence \(r_{10}=N_{10}=0\).
In this two-term situation \((r^{-1})_{10}=0\) as well,
since \(r(g^{-1})=r(g)^{-1}\). The symplectic adjoint formula
then gives \(r_{32}=0\). This excludes \(-A\) when \(B,P\)
are absent. The two exclusions leave only the first row
of the tight-root table.

\proofstep{Distinguish a global entry from an ordinary frame entry.}
The remaining off-diagonal degrees of the reduced matrices form a subset of
\[
 \{B,-P,-A,-T\}.
\] The frame has \(N_{21}=0\), since
its degree \(-B\) is absent. Tightness in degree \(B\)
comes from the global \((1,2)\) entry, since the frame
is lower triangular. The first lattice shape occurs
when degree \(-P\) appears in the global matrices.

In the second shape degree \(-P\) is tight only
through \(N_{20}\ne0\). Row two of the first-column
equation forces \(d_2=d_0\) on the marked decomposition
group. The global \((1,2)\) entry has the two-term
multiplication law. If it satisfied
\[
                      r_{12}=c(d_2-d_1)
\]
on this group for a scalar \(c\), normalized odd-root
vanishing would make it zero everywhere. Its nonzero
entry gives the required second lattice shape. \label{cor:cohomology-count-two-normalization-2}
\end{proof}
